\infographicchapter{Magnetic Particle Imaging (MPI)}{mpi.png}
\label{ch:infographic-mpi}

\section{Overview}

Magnetic particle imaging (MPI) is a new, experimental scanning method---invented in 2001 and first published in 2005---that creates three-dimensional images by tracking tiny superparamagnetic iron oxide nanoparticles (SPIONs) injected into the bloodstream. Unlike a regular MRI---which listens to signals from water inside your body---MPI detects only the injected iron particles. Because normal body tissues give off zero MPI signal, the iron tracer shines brightly against a completely dark, clutter-free background.

It is important to know that MPI is currently an investigational research technology and is not yet approved for routine medical diagnosis in hospitals or clinics. So far, scientists have tested MPI mainly in laboratory studies and early human research trials. In these studies, it shows exciting promise for fast, radiation-free blood vessel imaging, tracking stem cells or immune cells over days or weeks, finding hidden internal bleeding, locating lymph nodes, and even gently heating iron particles inside tumors.

MPI's biggest strengths are zero X-ray radiation, high sensitivity that is not blocked by bone or air, very fast image capture, and non-radioactive tracers that do not decay quickly. However, major hurdles still remain before routine hospital use: MPI does not show normal body anatomy on its own (so it must be paired with an MRI or CT scan), very few human-approved iron tracers exist, and building full human-sized scanners is technically challenging and expensive. In research studies, safety focuses on preventing mild nerve tingling or tissue warmth from changing magnetic fields, following strict MRI-style metal screening, and ensuring the iron particles are safe and cleared properly by the body.

\section{Underlying Physical Principles}

\begin{itemize}
    \item \textbf{Creating a Magnetic ``Sweet Spot'':} Strong magnets arranged opposite each other create a shaped magnetic field with a tiny zero-field spot (or line) in the middle---known technically as a \textit{field-free point} (FFP) or \textit{field-free line} (FFL). Iron nanoparticles outside this spot are locked in place magnetically, while particles inside the spot remain free to respond.
    \item \textbf{Flipping the Iron Particles:} A gentle, rapidly switching magnetic field (called the drive field) flips the magnetic direction of the free iron particles inside that central spot back and forth thousands of times per second. Normal body tissues do not respond to this flipping at all, leaving a completely clean background.
    \item \textbf{Picking Up the Unique Signal:} As the iron particles flip, they produce distinctive electrical ripples (harmonics) that sensitive wire coils pick up. By sweeping the magnetic sweet spot rapidly across the body, the scanner pinpoints exactly where the particles are located.
    \item \textbf{Measuring Exact Amounts:} The strength of the signal matches the exact amount of iron particles in each spot and passes straight through bone and air without weakening. A computer turns these signals into clear, 3D maps showing how much tracer is present at every point.
\end{itemize}

\section{Imaging Mechanism}

\begin{enumerate}
    \item \textbf{Scanner Positioning and Tracer Delivery:} In an approved research study, the participant lies inside the experimental MPI scanner (often paired with an MRI or CT scanner for anatomical reference), and a carefully measured dose of safe iron oxide nanoparticles is given through an arm IV or inside pre-labeled cells.
    \item \textbf{Sweeping the Magnetic Sweet Spot:} Strong coils create the shaped magnetic field with its central zero-field spot, while oscillating drive coils rapidly sweep that spot across the area being studied.
    \item \textbf{Listening for Particle Signals:} Sensitive receiver coils pick up the faint electrical signals created as the iron particles flip. Electronic filters remove background drive noise and convert the clean particle signals into digital data.
    \item \textbf{Building the 3D Tracer Map:} Reconstruction software translates the digital signals into a 3D map showing the exact concentration of iron particles in real time.
    \item \textbf{Overlaying on Anatomy:} Because MPI shows only the iron tracer and no surrounding organs, the bright color-coded MPI map is lined up and layered on top of a matching MRI or CT scan so researchers can see exactly where the tracer sits in the body.
\end{enumerate}

\section{Key Technical Parameters}

\noindent\begin{tabularx}{\textwidth}{@{}l X@{}}
\toprule
\textbf{Parameter} & \textbf{Typical Value / Description} \\
\midrule
Physical Principle & Detects the magnetic flipping signal of injected iron oxide nanoparticles as a zero-field magnetic spot sweeps through the body \\
Energy / Signal Source & Strong static magnetic gradient field paired with rapidly oscillating magnetic drive fields \\
Typical Frequency / Energy & Oscillating magnetic fields at 1--25 kilohertz (kHz); no X-rays or radioisotopes \\
Spatial Resolution & About 0.5--2.0 mm in preclinical research scanners (depends on magnet strength and particle design) \\
Typical Scan Duration & Seconds to a few minutes per MPI scan (30--60 minutes for a full research visit with paired MRI or CT) \\
Contrast Agent & Required iron oxide nanoparticle tracer (given by IV or attached to tracked cells) \\
Ionizing Radiation & None (zero X-ray or radioactive exposure) \\
\bottomrule
\end{tabularx}

\section{Clinical Usages}

\subsection{Primary Usages}

\begin{itemize}
    \item \textbf{Blood Vessel Imaging (Investigational):} MPI is not yet approved for routine medical care. One of its top research goals is fast, radiation-free 3D imaging of blood vessels and narrowed arteries using iron particles circulating in the bloodstream.
    \item \textbf{Brain and Heart Blood Flow (Investigational):} Preclinical studies are testing MPI to spot sudden stroke, reduced brain blood flow, or heart blood-flow problems very quickly, without the skull or ribs blocking the signal. Academic medical centers are also building human-scale brain MPI scanners designed to detect cerebral blood-volume changes with roughly tenfold higher sensitivity than conventional functional MRI (at about 5 mm spatial resolution and 5-second speed) to help assess stroke size, monitor brain bleeding, and map individual brain activity.
    \item \textbf{Tracking Therapeutic Cells (Investigational):} Researchers use MPI to follow iron-labeled stem cells, CAR-T cancer-fighting cells, and immune cells inside the body over days or weeks without using radioactive tracers that fade away.
    \item \textbf{Spotting Internal Bleeding and Lung Clots (Investigational):} Early lab studies show high sensitivity for catching active stomach or intestinal bleeding (as tracer leaks out of a vessel) and blood-flow gaps in the lungs without X-ray radiation.
\end{itemize}

\subsection{Secondary Usages}

\begin{itemize}
    \item \textbf{Lymph Node and Tumor Research:} Preclinical studies explore using MPI to locate sentinel lymph nodes and see how iron particles collect inside solid tumors.
    \item \textbf{Checking Tumor Margins During Surgery (Investigational):} Prototype MPI tools are being studied for use during breast-conserving cancer surgery (lumpectomy) to quickly check whether tumor cells remain at the edges of removed tissue, with the goal of lowering the need for repeat operations.
    \item \textbf{Targeted Tumor Heating (Theranostics):} Scientists are studying whether the same magnetic sweet spot used for imaging can also gently heat iron particles right inside a tumor to destroy cancer cells while sparing healthy nearby tissue.
    \item \textbf{Tracking Medical Wires and Catheters:} Research teams are testing iron-coated guidewires and catheters so doctors could one day steer vascular tools in real time without X-ray exposure.
    \item \textbf{Hybrid Research Scans (MPI/MRI and MPI/CT):} Combining MPI with MRI or CT in research scanners pairs exact iron-tracer measurements with detailed maps of body organs.
\end{itemize}

\section{Advantages}

\begin{itemize}
    \item \textbf{Zero Radiation and Clean Background:} Uses magnetic fields and iron particles instead of X-rays or radioactivity, producing crisp, glowing 3D tracer images with zero distracting signal from normal tissues.
    \item \textbf{Exact Measurement at Any Depth:} The brightness of the image directly matches the amount of iron tracer present and is not blocked by thick bone, air-filled lungs, or deep tissue.
    \item \textbf{High Sensitivity and Real-Time Speed:} Detects tiny amounts of iron nanoparticles and captures images in fractions of a second, making it promising for fast blood-flow and moving-device studies.
    \item \textbf{Long-Lasting, Non-Radioactive Tracers:} Unlike nuclear medicine tracers that decay within hours, iron oxide nanoparticles stay stable on the shelf and allow cell tracking over days or weeks.
    \item \textbf{Combined Imaging and Treatment Potential:} In research settings, the same scanner hardware can both locate iron particles and gently heat them at a precise target spot.
\end{itemize}

\section{Disadvantages}

\begin{itemize}
    \item \textbf{Experimental Status Only:} MPI is still in preclinical and early human research stages and is not approved by regulators for routine clinical diagnosis.
    \item \textbf{Shows No Normal Anatomy on Its Own:} Because normal organs give off no signal, an MPI scan cannot show bones or soft tissues by itself and must be paired with an MRI or CT scan.
    \item \textbf{Hard to Scale to Full Human Size:} Achieving sharp resolution requires very strong magnetic gradients, making human-sized scanners complex, power-hungry, and costly to build.
    \item \textbf{Very Few Approved Tracers:} Most available iron nanoparticle agents were designed for MRI rather than MPI, and new MPI-specific tracers still need thorough safety testing for allergy risk and iron clearance.
    \item \textbf{Nerve Tingling, Heating, and Magnet Limits:} Rapidly switching magnetic fields can cause mild nerve tingling or tissue warming if pushed too high, and strong magnets bring the same strict metallic implant rules as MRI.
\end{itemize}

\section{Contraindications}

\subsection{Absolute}

\begin{itemize}
    \item \textbf{Outside Approved Research, Unsafe Metal Implants, or Severe Allergy:} MPI must not be used for routine clinical care outside approved research protocols (IRB and FDA Investigational Device Exemption). Standard MRI-style metal bans apply---including non-conditional pacemakers or defibrillators, magnetic brain aneurysm clips, and metal fragments in the eyes---as well as any severe past allergy to iron oxide nanoparticle agents.
\end{itemize}

\subsection{Relative}

\begin{itemize}
    \item \textbf{Iron Overload, Liver Disease, Asthma, and Pregnancy:} Conditions with too much iron in the body (such as hemochromatosis), severe liver disease (which slows clearance of iron particles), severe uncontrolled asthma, and pregnancy or breastfeeding are generally excluded from human MPI research trials.
    \item \textbf{Comfort Limits:} Severe claustrophobia or low tolerance for mild tingling or muscle twitches from switching magnetic fields.
\end{itemize}

\section{Patient Safety Guidelines}

\begin{itemize}
    \item \textbf{Strict Research Oversight:} Because MPI is investigational, any human scan must take place under an Institutional Review Board (IRB)-approved study with regulatory oversight (such as an FDA Investigational Device Exemption) and full informed consent.
    \item \textbf{Preventing Nerve Tingling and Tissue Heating:} Scanners are programmed to stay within strict magnetic-field and heating limits so switching fields do not cause painful nerve stimulation (tingling or muscle twitches) or excess warmth.
    \item \textbf{Safe Iron Tracer Handling:} Only carefully tested iron nanoparticle tracers are given slowly under close medical watch. Staff screen for iron allergies, asthma, and iron-overload conditions, and keep emergency allergy medicines ready.
    \item \textbf{MRI-Style Metal Screening:} Strict magnet safety zones and screening checklists keep loose metal objects out of the room and protect anyone with pacemakers, clips, or other metal implants.
    \item \textbf{Close Monitoring and Anatomical Pairing:} Research teams monitor vital signs continuously, keep quick-exit procedures ready, and always pair MPI data with standard MRI or CT scans before drawing any research conclusions.
\end{itemize}

\section{Observing for Abnormal Patterns}

\begin{itemize}
    \item \textbf{Narrowed Vessels and Blood-Flow Gaps (Research Findings):} In lab studies of narrowed arteries, stroke, or lung clots, a sudden pinching of the bright blood-vessel line points to narrowing, while a sharp dark gap shows blocked blood flow---clearly visible even through the skull or ribs.
    \item \textbf{Active Internal Bleeding:} In preclinical bleeding studies, a bright glowing spot appearing outside a blood vessel and slowly growing over time pinpoints active internal bleeding with very high sensitivity.
    \item \textbf{Cell Tracking and Tumor Uptake:} Glowing hotspots that linger hours or days later show where iron-labeled stem cells or immune cells have traveled, or where particles have settled inside inflamed areas, lymph nodes, or tumors.
    \item \textbf{Normal Liver and Spleen Glow:} Because the liver and spleen normally collect iron particles as they clear them from the blood, bright signal there is expected and can sometimes hide nearby upper-belly targets unless carefully lined up with MRI or CT.
    \item \textbf{Technical Look-Alikes:} Slight alignment shifts between MPI and MRI/CT can make a glowing spot look like it sits in a neighboring organ. Clumping of nanoparticles can blur the image slightly, and the very bright liver signal can overshadow fainter signals nearby.
\end{itemize}

\section{Dealing with Uncertainties}

\begin{itemize}
    \item \textbf{Not Yet Standardized for Medical Decisions:} Because MPI is still experimental, scanning settings and interpretation rules are not yet standardized for routine patient care, so MPI results alone are never used to make clinical treatment decisions.
    \item \textbf{Alignment and Background Challenges:} Because MPI does not show organs directly, even a tiny shift when overlaying it on MRI or CT---or bright normal uptake in the liver and spleen---can make exact localization tricky.
    \item \textbf{Particle Behavior Variations:} If nanoparticles clump together or clear at different speeds, the signal can blur slightly or become harder to measure accurately.
    \item \textbf{Confirming with Approved Tests:} Research teams double-check MPI scans using calibration phantoms and careful MRI/CT alignment, and any suspected health issue found during a study is always confirmed with approved standard tests like CT, MRI, ultrasound, or biopsy.
\end{itemize}

\section{Augmenting Methods}

\begin{itemize}
    \item \textbf{Paired MRI and CT Scans:} Lining up MPI with high-detail MRI or CT is essential so researchers can see exactly which organs and blood vessels hold the glowing iron tracer.
    \item \textbf{Standard Vascular and Nuclear Scans:} Established clinical tests like CT angiography, MR angiography, Doppler ultrasound, and PET or SPECT serve as trusted benchmarks to validate investigational MPI results.
    \item \textbf{``Multi-Color'' MPI Processing:} Advanced research software can analyze subtle differences in how fast nanoparticles relax, helping separate free-flowing particles from bound ones or even estimate local temperature during experimental heating studies.
    \item \textbf{Lab Checks and Tissue Validation:} In research studies, iron blood tests, liver and kidney panels, and microscopic tissue checks verify how safely and completely the nanoparticles are cleared.
\end{itemize}

\section{Comparison with Alternative Modalities}

\begin{itemize}
    \item \textbf{Compared with MRI:} MRI is proven worldwide and shows rich anatomical detail of soft tissues. MPI cannot show normal anatomy on its own, but it tracks iron nanoparticles with much higher sensitivity, exact measurement, and no dark distortion artifacts.
    \item \textbf{Compared with PET and SPECT:} Nuclear scans like PET and SPECT are clinically proven for molecular imaging but use ionizing radiation and short-lived radioactive tracers. MPI aims for similar sensitivity using shelf-stable, non-radioactive iron particles at much faster speeds.
    \item \textbf{Compared with CT Angiography:} CT angiography is the fast, widely available standard for checking blood vessels today. Investigational MPI aims to eventually offer radiation-free, iodine-free, real-time 3D blood-vessel imaging that is not blocked by bone.
\end{itemize}

\section{Key References}

\begin{itemize}
    \item Gleich B, Weizenecker J. Tomographic Imaging Using the Nonlinear Response of Magnetic Particles. \textit{Nature}, 2005;435(7046):1214--1217.
    \item Saritas EU, Goodwill PW, Croft LR, et al. Magnetic Particle Imaging (MPI) for NMR and MRI Researchers. \textit{Journal of Magnetic Resonance}, 2013;229:116--126.
    \item Billings C, Langley M, Warrington G, Mashali F, Johnson JA. Magnetic Particle Imaging: Current and Future Applications, Magnetic Nanoparticle Synthesis Methods and Safety Measures. \textit{International Journal of Molecular Sciences}, 2021;22(14):7651.
    \item Graeser M, Thieben F, Szwargulski P, et al. Human-Sized Magnetic Particle Imaging for Brain Applications. \textit{Nature Communications}, 2019;10:1936.
    \item Mass General Advances in Motion. \textit{Magnetic Particle Imaging: New Imaging Modality Could Advance Stroke Assessment, Other Applications}. Massachusetts General Hospital. Available at: \url{https://advances.massgeneral.org/radiology/article.aspx?id=1555}.
\end{itemize}

\section{Further Reading}
\begin{itemize}
  \item Wikipedia contributors. \textit{Magnetic particle imaging}. Wikipedia, The Free Encyclopedia. Available at: \url{https://en.wikipedia.org/wiki/Magnetic_particle_imaging}.
\end{itemize}
